\documentclass[12pt,a4paper]{article}

% Packages
\usepackage[utf8]{inputenc}
\usepackage[T1]{fontenc}
\usepackage{lmodern}
\usepackage[margin=2.5cm]{geometry}
\usepackage{amsmath,amssymb}
\usepackage{graphicx}
\usepackage{booktabs}
\usepackage{natbib}
\usepackage{hyperref}
\usepackage{setspace}
\usepackage{enumitem}
\usepackage{longtable}
\usepackage{float}

% Document settings
\onehalfspacing
\setlength{\parindent}{0pt}
\setlength{\parskip}{0.5em}

% Title
\title{\textbf{Structural Barriers to Scientific Integration:\\A Metascientific Analysis of Behavioral Economics Fragmentation}}
\author{Evidence-Based Framework Research Team\thanks{FehrAdvice \& Partners AG. Correspondence: research@fehradvice.com}\\[0.5em]\small Working Paper WP-2026-02 (Full Version)}
\date{January 2026}

\begin{document}

\maketitle

\begin{abstract}
Why do 40+ behavioral models exist in isolation despite universal recognition that integration would benefit science and practice? This paper argues that the barrier is structural, not intellectual. Drawing on metascience, sociology of knowledge, philosophy of science, and science of science literatures, we identify five structural barriers: (1) academic incentive structures favoring novelty over synthesis, (2) disciplinary silos with incompatible formalisms, (3) systematic peer review bias against interdisciplinary work, (4) funding structures rewarding narrow hypotheses, and (5) territorial defense of paradigmatic authority. We document evidence from citation network analysis, publication patterns, experimental studies of peer review, and historical case studies.

A central contribution is the \textit{framework versus model distinction}: classical economics (Samuelson 1947) was conceived as a framework but treated as a model, leading to fragmentation when anomalies accumulated. We trace this epistemological confusion through the Chicago School's methodological development and show why it was technologically defensible until recent computational advances made non-separable, context-dependent modeling feasible.

We conclude that integration requires institutional innovation outside traditional academic structures---a prediction consistent with the emergence of integrative frameworks in consulting and policy contexts rather than universities. The Evidence-Based Framework for Economic and Social Behavior (EBF) exemplifies this pattern, built outside academia with different incentive structures and validated through practice rather than publication.

\medskip
\textbf{Keywords:} metascience, scientific integration, behavioral economics, incentive structures, disciplinary silos, framework versus model, Chicago School, replication crisis

\medskip
\textbf{JEL Codes:} A11, A14, B21, B41, D91
\end{abstract}

\newpage
\tableofcontents
\newpage

%==============================================================================
\section{Introduction: The Integration Puzzle}
\label{sec:intro}
%==============================================================================

\subsection{The Puzzle}

Behavioral science has accumulated remarkable knowledge about human decision-making over the past half-century:

\begin{itemize}
\item \textbf{Psychology:} Transtheoretical Model, Health Action Process Approach, Rubicon Model, Self-Determination Theory, Protection Motivation Theory (6+ stage models of behavior change)
\item \textbf{Economics:} Becker-Murphy rational addiction, Fudenberg-Levine dual-self, HANK heterogeneous agents, Gabaix inattention models (10+ formal models)
\item \textbf{Sociology:} Granovetter threshold models, information cascades, SIR contagion dynamics
\item \textbf{Management:} Lewin change models, Schein organizational culture, PARC, COM-B frameworks
\item \textbf{Neuroscience:} Dual-process theories, neural autopilot, habit formation mechanisms
\end{itemize}

Yet these models remain \textbf{fragmented}:
\begin{itemize}
\item Different disciplines use different terminology for identical concepts
\item Mathematical formalisms are incompatible across fields
\item No cumulative validation exists across frameworks
\item Practitioners receive contradictory advice from different experts
\end{itemize}

\subsection{The Core Paradox}

\begin{quote}
\textit{Everyone agrees integration is needed. No one is rewarded for doing it.}
\end{quote}

This fragmentation persists despite widespread agreement that integration would be valuable. The puzzle is not whether integration is desirable---it clearly is---but why it has not occurred.

\subsection{The Answer: Structural Barriers}

This paper argues that the lack of integration is \textbf{not} due to:
\begin{itemize}
\item Insufficient knowledge (we have enough models)
\item Technical impossibility (the existence of integrative frameworks proves it's possible)
\item Lack of demand (practitioners desperately need unified guidance)
\end{itemize}

Rather, it is due to \textbf{structural barriers} in the academic system that systematically discourage integration while rewarding fragmentation.

\subsection{The Hypothesis}

We propose that the barrier to integration is \textit{structural}, not intellectual:

\begin{quote}
\textit{Academic institutions systematically reward fragmentation and penalize integration through incentive structures, disciplinary boundaries, peer review processes, funding mechanisms, and territorial dynamics.}
\end{quote}

This hypothesis has a testable implication: if the barrier is structural rather than intellectual, integration should emerge from contexts with different incentive structures---consulting firms, policy organizations, or research institutes outside traditional academia.

\subsection{Contribution}

This paper synthesizes evidence from five literatures:
\begin{enumerate}[nosep]
\item Philosophy of science (Kuhn, Popper, Lakatos)
\item Sociology of scientific knowledge (Merton, Bloor, Latour)
\item Economics of science (Stephan, Smaldino \& McElreath)
\item Metascience (Ioannidis, Nosek, Open Science Collaboration)
\item Science of science (Fortunato et al., Uzzi et al., Wang et al.)
\end{enumerate}

We document five structural barriers, present quantitative evidence for each, and trace the historical epistemological confusion that contributed to fragmentation.

%==============================================================================
\section{Classical Foundations: Kuhn, Popper, and Lakatos}
\label{sec:classics}
%==============================================================================

\subsection{Kuhn: Paradigms and Normal Science (1962)}

Thomas Kuhn's \textit{The Structure of Scientific Revolutions} provides the foundational framework for understanding why science fragments.

\textbf{Key Concepts:}
\begin{itemize}
\item \textbf{Paradigm:} A shared framework of concepts, methods, and exemplars
\item \textbf{Normal Science:} Puzzle-solving within the paradigm
\item \textbf{Incommensurability:} Different paradigms cannot be directly compared
\end{itemize}

\textbf{Implication for Integration:}
\begin{quote}
``Normal science... is predicated on the assumption that the scientific community knows what the world is like. Much of the success of the enterprise derives from the community's willingness to defend that assumption.'' (Kuhn, 1962, p. 5)
\end{quote}

Kuhn explains why \textit{within}-paradigm work is rewarded: it solves recognized puzzles. \textit{Between}-paradigm work is seen as threatening, not contributing.

\subsection{Popper: Falsification and Demarcation (1959)}

Karl Popper's falsificationism emphasizes testable hypotheses:
\begin{itemize}
\item Science progresses through conjecture and refutation
\item Good theories make bold, falsifiable predictions
\item Unfalsifiable theories are ``not even wrong''
\end{itemize}

\textbf{Problem for Integration:}
Meta-frameworks are harder to falsify than narrow hypotheses. A claim like ``this framework unifies 40 models'' is harder to test than ``Effect X = 0.3 in population P.''

The Popperian criterion inadvertently favors fragmentation by rewarding narrow, testable claims over broad integrative ones.

\subsection{Lakatos: Research Programmes (1970)}

Imre Lakatos attempted to reconcile Kuhn and Popper:
\begin{itemize}
\item \textbf{Hard Core:} Unfalsifiable central assumptions
\item \textbf{Protective Belt:} Auxiliary hypotheses that can be modified
\item \textbf{Progressive vs. Degenerating:} Programmes that predict novel facts vs. those that don't
\end{itemize}

\textbf{Application to Integration:}
An integrative framework can be seen as a new research programme with:
\begin{itemize}
\item Hard Core: Central organizing concepts (e.g., complementarity, context-dependence)
\item Protective Belt: Specific parameter estimates for each context
\item Progressive: Predicts novel facts that component models cannot predict alone
\end{itemize}

%==============================================================================
\section{Sociology of Scientific Knowledge}
\label{sec:ssk}
%==============================================================================

\subsection{Merton: The Normative Structure of Science (1942)}

Robert Merton identified four norms of science (CUDOS):
\begin{enumerate}
\item \textbf{Communism:} Knowledge is shared
\item \textbf{Universalism:} Claims evaluated on merit, not source
\item \textbf{Disinterestedness:} Scientists pursue truth, not self-interest
\item \textbf{Organized Skepticism:} Claims are scrutinized
\end{enumerate}

\textbf{Reality Check:}
Modern science violates all four:
\begin{itemize}
\item Communism $\to$ Proprietary data, paywalled journals
\item Universalism $\to$ Prestige bias, Matthew Effect
\item Disinterestedness $\to$ Career incentives dominate
\item Organized Skepticism $\to$ Peer review is tribal
\end{itemize}

\subsection{The Strong Programme: Bloor (1976)}

David Bloor's ``strong programme'' in sociology of science argues:
\begin{itemize}
\item Scientific knowledge is socially constructed
\item The same sociological methods explain ``true'' and ``false'' beliefs
\item Science is not exempt from sociological analysis
\end{itemize}

\textbf{Implication:}
The fragmentation of behavioral science is not a failure of logic but a \textit{social} outcome of institutional structures.

\subsection{Latour and Woolgar: Laboratory Life (1979)}

Bruno Latour and Steve Woolgar's ethnography of a biology lab showed:
\begin{itemize}
\item ``Facts'' are constructed through social processes
\item Inscription devices (instruments, papers) create ``black boxes''
\item Once a claim becomes a ``fact,'' its social origins are erased
\end{itemize}

\textbf{For Integration:}
Each behavioral model is a ``black box'' whose social origins (disciplinary context, career incentives, paradigmatic commitments) are invisible. Integration requires opening these boxes.

%==============================================================================
\section{Five Structural Barriers to Integration}
\label{sec:barriers}
%==============================================================================

\subsection{Theoretical Framework}

Define scientific integration as the construction of a meta-framework $\mathcal{F}$ that:
\begin{enumerate}[nosep]
\item Subsumes existing models $\{M_1, M_2, \ldots, M_n\}$ as special cases
\item Uses consistent notation and assumptions
\item Enables cumulative validation across domains
\item Generates novel predictions beyond component models
\end{enumerate}

Integration is a public good: costly to produce, freely available once produced, and beneficial to all users. Standard public goods theory predicts underproduction.

\subsection{A Simple Model}

Let $V_I$ denote the social value of integration and $C_I$ its cost to the integrating scientist. Let $p$ denote the probability of successful publication and career advancement.

Integration occurs if:
\begin{equation}
p \cdot V_{\text{private}} > C_I
\end{equation}

where $V_{\text{private}} \ll V_I$ because most benefits of integration are external.

Under current academic incentives:
\begin{itemize}[nosep]
\item $p$ is low (peer review bias, disciplinary gatekeeping)
\item $V_{\text{private}}$ is low (integration papers don't fit career metrics)
\item $C_I$ is high (requires expertise across multiple fields)
\end{itemize}

The model predicts severe underproduction of integration.

%------------------------------------------------------------------------------
\subsection{Barrier 1: Academic Incentive Structures}
\label{sec:barrier1}
%------------------------------------------------------------------------------

\begin{equation}
\text{Academic Success} = f(\text{Publications}, \text{Citations}, \text{Grants})
\end{equation}

None of these metrics directly reward integration.

\subsubsection{Smaldino \& McElreath (2016): Natural Selection of Bad Science}

Paul Smaldino and Richard McElreath developed an evolutionary model of science:

\textbf{Model:}
\begin{itemize}
\item Labs ``reproduce'' (train students) based on publication success
\item Selection pressure favors methods that produce more publications
\item Quality (replicability, integration) is not directly selected for
\end{itemize}

\textbf{Key Finding:}
\begin{quote}
``As long as the incentives are in place that reward publication quantity... we expect that the rate of firing [researchers using poor methods] would have to be implausibly high to prevent the system from selecting for poor methods.''
\end{quote}

\textbf{Implications:}
\begin{enumerate}
\item Bad methods persist because they are \textit{productive}
\item Good methods (like integration) are selected \textit{against}
\item The system is in a stable equilibrium of fragmentation
\end{enumerate}

\subsubsection{Stephan (2012): How Economics Shapes Science}

Paula Stephan's comprehensive analysis shows:
\begin{itemize}
\item Scientists respond rationally to incentives
\item Incentives favor safe, incremental work
\item High-risk integration is career-threatening
\item The ``winner-take-all'' job market increases risk aversion
\end{itemize}

\subsubsection{Edwards \& Roy (2017): Perverse Incentives}

\begin{table}[H]
\centering
\caption{Perverse Incentives in Academia (Edwards \& Roy, 2017)}
\label{tab:perverse}
\renewcommand{\arraystretch}{1.3}
\begin{tabular}{@{}ll@{}}
\toprule
\textbf{Metric} & \textbf{Perverse Effect} \\
\midrule
Publication count & Salami-slicing, least publishable units \\
Journal impact factor & Chasing ``hot'' topics, avoiding integration \\
Citation count & Self-citation, citation cartels \\
Grant funding & Safe projects, avoiding risk \\
H-index & Gaming, quantity over quality \\
\bottomrule
\end{tabular}
\end{table}

%------------------------------------------------------------------------------
\subsection{Barrier 2: Disciplinary Silos}
\label{sec:barrier2}
%------------------------------------------------------------------------------

\subsubsection{The Structure of Silos}

Disciplines develop specialized vocabularies, methods, and standards of evidence (Klein, 1990). These serve legitimate functions---enabling cumulative progress within paradigms---but create barriers to cross-disciplinary communication.

The same concept receives different names in different fields:
\begin{itemize}
\item ``Habit'' (psychology) = ``Autopilot'' (neuroscience) = ``State-dependence'' (economics)
\item ``Social proof'' (psychology) = ``Information cascade'' (economics) = ``Conformity'' (sociology)
\item ``Ego depletion'' (psychology) = ``Self-control depletion'' (economics) = ``Decision fatigue'' (management)
\end{itemize}

\subsubsection{Jacobs (2013): In Defense of Disciplines}

Jerry Jacobs argues that disciplines exist for good reasons:
\begin{itemize}
\item Shared methods enable cumulative progress
\item Specialized training ensures expertise
\item Peer review works within shared paradigms
\end{itemize}

But he also notes:
\begin{quote}
``The very mechanisms that make disciplines productive also create barriers to the flow of ideas across fields.''
\end{quote}

\subsubsection{Klein (1990): Interdisciplinarity}

Julie Thompson Klein's comprehensive analysis identifies barriers:
\begin{itemize}
\item \textbf{Cognitive:} Different assumptions, methods, standards of evidence
\item \textbf{Social:} Different communities, conferences, journals
\item \textbf{Institutional:} Different departments, hiring criteria, tenure standards
\item \textbf{Political:} Competition for resources, status hierarchies
\end{itemize}

\subsubsection{Fortunato et al. (2018): Science of Science}

Network analysis of 21 million papers shows:
\begin{itemize}
\item Citation networks cluster strongly by discipline
\item Cross-disciplinary citations are rare and declining
\item ``Star'' scientists bridge silos but are exceptions
\end{itemize}

Formally:
\begin{equation}
P(\text{cite across discipline}) \approx 0.1 \cdot P(\text{cite within discipline})
\end{equation}

%------------------------------------------------------------------------------
\subsection{Barrier 3: Peer Review Bias Against Integration}
\label{sec:barrier3}
%------------------------------------------------------------------------------

\subsubsection{Uzzi et al. (2013): Atypical Combinations}

Brian Uzzi and colleagues analyzed 17.9 million papers:

\textbf{Finding 1:} High-impact papers combine conventional and atypical knowledge.

\textbf{Finding 2:} Reviewers systematically penalize atypicality.

\begin{quote}
``Papers that are highly novel are more likely to be rejected, take longer to get published, and are initially cited less—but eventually have higher impact if they survive.''
\end{quote}

\textbf{Implication:} Integration is high-risk, high-reward—but the risk is borne by the researcher, the reward accrues to society.

\subsubsection{Wang et al. (2017): Bias Against Novelty}

Jian Wang and colleagues document quantitative bias against novelty:
\begin{itemize}
\item Novel papers take 15\% longer to publish
\item Novel papers are 20\% more likely to be desk-rejected
\item Novel papers have higher variance in outcomes
\item Effect sizes are larger for interdisciplinary work
\end{itemize}

\subsubsection{The Reviewer's Dilemma}

\textbf{Typical reviewer reasoning for integration papers:}
\begin{itemize}
\item ``This is too broad---not rigorous enough for our discipline''
\item ``The author doesn't understand [my field] deeply enough''
\item ``This doesn't fit our journal's scope''
\item ``Where is the novel empirical contribution?''
\end{itemize}

\textbf{Result:} Integration papers are rejected by \textit{all} disciplines because they fail to meet each discipline's specialized criteria.

%------------------------------------------------------------------------------
\subsection{Barrier 4: Funding Structures}
\label{sec:barrier4}
%------------------------------------------------------------------------------

Grant agencies reward well-defined, testable hypotheses. Integration is inherently risky: outcomes are uncertain, timelines are long, and success criteria are ambiguous. Risk-averse funding committees favor incremental projects.

Analysis of NIH and NSF funding patterns reveals:
\begin{itemize}
\item Success rates decline for interdisciplinary proposals
\item Review panels lack expertise to evaluate integration
\item Safe, incremental projects dominate portfolios
\item Integration requires longer time horizons than typical 3-year grants
\end{itemize}

Junior researchers---who might have the energy for integration---face the strongest pressure to conform because the academic job market is winner-take-all.

%------------------------------------------------------------------------------
\subsection{Barrier 5: Territorial Defense}
\label{sec:barrier5}
%------------------------------------------------------------------------------

Paradigms are not merely intellectual structures but social institutions with defenders. Integration threatens established authority by suggesting that existing models are incomplete. Kuhn (1962) documented how paradigm defenders resist challenges.

Historical case studies document paradigm defense:
\begin{itemize}
\item \textbf{Behavioral economics} faced decades of resistance from rational choice theorists who controlled top journals
\item \textbf{Plate tectonics} was rejected for 50 years despite geological evidence
\item \textbf{Ulcer-bacteria theory} met fierce opposition from the acid-reduction paradigm
\item \textbf{Prospect theory} took years to be accepted despite experimental evidence
\end{itemize}

The pattern is consistent: integration threatens established authority and faces resistance proportional to the stakes involved.

%==============================================================================
\section{The Replication Crisis as Symptom}
\label{sec:replication}
%==============================================================================

\subsection{Ioannidis (2005): Why Most Research Findings Are False}

John Ioannidis's landmark paper showed:
\begin{itemize}
\item For most study designs, most published findings are false
\item Small sample sizes, small effect sizes, flexibility in analysis
\item Publication bias: Positive results published, null results buried
\end{itemize}

\textbf{Connection to Fragmentation:}
Each fragmented model is validated only within its silo. Without integration, there is no cumulative validation across frameworks.

\subsection{Open Science Collaboration (2015)}

The Reproducibility Project attempted to replicate 100 psychology studies:
\begin{itemize}
\item Only 36\% replicated successfully
\item Mean effect size dropped by 50\%
\item High-profile findings often failed
\end{itemize}

\subsection{Camerer et al. (2016, 2018): Economics Replications}

Similar findings in economics:
\begin{itemize}
\item 61\% of experimental economics studies replicated (2016)
\item 62\% of social science papers in Nature/Science replicated (2018)
\end{itemize}

\textbf{Implication:}
Fragmented models cannot be cumulatively validated. Integration would enable cross-model consistency checks that might catch replication failures earlier.

%==============================================================================
\section{Scholarly Calls for Integration}
\label{sec:calls}
%==============================================================================

\subsection{Wilson (1998): Consilience}

Edward O. Wilson's manifesto for unification:
\begin{quote}
``The greatest enterprise of the mind has always been and always will be the attempted linkage of the sciences and humanities. The ongoing fragmentation of knowledge and resulting chaos in philosophy are not reflections of the real world but artifacts of scholarship.''
\end{quote}

Wilson argues that biology, psychology, and social science should converge---they study the same phenomenon (human behavior) at different levels.

\subsection{Gintis (2009): The Bounds of Reason}

Herbert Gintis explicitly calls for behavioral science unification:
\begin{quote}
``The behavioral sciences should be unified... The current separation of economics, psychology, sociology, and anthropology is an artifact of disciplinary history, not a reflection of natural boundaries in the subject matter.''
\end{quote}

Gintis proposes:
\begin{itemize}
\item Shared mathematical framework (game theory + evolution)
\item Common experimental methods
\item Unified theory of human behavior
\end{itemize}

\subsection{Bammer (2013): Integration and Implementation Sciences}

Gabriele Bammer proposes a new field dedicated to integration:
\begin{itemize}
\item \textbf{Integration:} Combining knowledge from different disciplines
\item \textbf{Implementation:} Applying integrated knowledge to real problems
\item \textbf{I2S:} A discipline whose subject matter is integration itself
\end{itemize}

\subsection{Other Calls}

\begin{table}[H]
\centering
\caption{Scholarly Calls for Integration}
\label{tab:calls}
\renewcommand{\arraystretch}{1.3}
\begin{tabular}{@{}lll@{}}
\toprule
\textbf{Author} & \textbf{Work} & \textbf{Call} \\
\midrule
Wilson (1998) & Consilience & Unity of knowledge \\
Gintis (2009) & Bounds of Reason & Unified behavioral science \\
Bammer (2013) & Disciplining Interdisciplinarity & Integration as discipline \\
Arthur (2015) & Complexity Economics & Complexity as unifier \\
D.S. Wilson (2019) & This View of Life & Evolution as meta-theory \\
\bottomrule
\end{tabular}
\end{table}

%==============================================================================
\section{The Framework versus Model Distinction}
\label{sec:framework-model}
%==============================================================================

\subsection{The Epistemological Core}

Before examining how to address structural barriers, we must clarify a fundamental epistemological distinction that explains \emph{why} economics fragmented.

\begin{table}[H]
\centering
\caption{Framework versus Model}
\label{tab:framework-model}
\renewcommand{\arraystretch}{1.4}
\begin{tabular}{@{}p{3cm}p{5.5cm}p{5.5cm}@{}}
\toprule
& \textbf{Framework} & \textbf{Model} \\
\midrule
\textbf{Purpose} & Organizes thinking & Makes predictions \\
\textbf{Claim} & ``Here is how to think about X'' & ``X equals Y'' \\
\textbf{Falsifiable?} & No (meta-level) & Yes (testable) \\
\textbf{On anomalies} & Integrates them & Is falsified by them \\
\textbf{Scope} & Broad, organizing & Narrow, specific \\
\textbf{Example} & Newtonian mechanics as worldview & $F = ma$ as equation \\
\bottomrule
\end{tabular}
\end{table}

\subsection{The Map Analogy}

\textbf{A Framework is like a map style} (topographic, political, road map).

\textbf{A Model is like a specific route} (``Take A1 to Bern, exit 24'').

\begin{table}[H]
\centering
\renewcommand{\arraystretch}{1.3}
\begin{tabular}{@{}p{6cm}p{6cm}@{}}
\toprule
\textbf{Map Style (Framework)} & \textbf{Route (Model)} \\
\midrule
Shows \textit{what to look for} & Tells \textit{what to do} \\
Multiple routes possible & One specific prediction \\
Can't be ``wrong''---only useful or not & Can be wrong (road closed) \\
Integrates new roads easily & Fails if road doesn't exist \\
\bottomrule
\end{tabular}
\end{table}

\textbf{Key insight:} When your route fails (traffic jam), you don't throw away the map. You use the map to find another route. But economics threw away the map when routes failed.

\subsection{The Historical Confusion: Samuelson 1947}

Paul Samuelson's \emph{Foundations of Economic Analysis} (1947) was conceived as a \textbf{framework}---a mathematical language for expressing economic relationships. Samuelson himself noted the generality of the approach.

However, the profession \textbf{treated it as a model}:
\begin{itemize}[nosep]
\item $C_{ij} = 0$ (separability) became a \emph{testable assumption}
\item Rational utility maximization became a \emph{prediction}
\item Equilibrium existence became an \emph{empirical claim}
\end{itemize}

\subsection{The Consequence: Fragmentation}

When anomalies accumulated (Allais 1953, Kahneman-Tversky 1979, G\"{u}th 1982, Fehr-Schmidt 1999):
\begin{enumerate}[nosep]
\item The ``model'' was falsified by experimental evidence
\item But no \textbf{framework} existed to integrate the anomalies
\item Each anomaly became its own literature (behavioral economics, identity economics, institutional economics...)
\item Result: \textbf{Fragmentation}
\end{enumerate}

Classical economics was a \textbf{framework} conceived as a way of thinking. It was \textbf{treated as a model} with testable predictions. When predictions failed, \textbf{it could serve as neither}:
\begin{itemize}[nosep]
\item Not as a model (falsified)
\item Not as a framework (couldn't integrate anomalies)
\end{itemize}

Result: Economics lost its \textbf{gravitational center}.

%==============================================================================
\section{The Chicago School: Documentation of the Framework-Model Confusion}
\label{sec:chicago}
%==============================================================================

\subsection{The Chicago School Generations}

The ``Chicago School'' was not monolithic but evolved through distinct generations, each reinforcing the treatment of separability ($C_{ij} = 0$) as empirical fact rather than simplifying assumption.

\begin{table}[H]
\centering
\caption{Chicago School Generations and Key Figures}
\label{tab:chicago-generations}
\renewcommand{\arraystretch}{1.3}
\small
\begin{tabular}{@{}llp{6cm}@{}}
\toprule
\textbf{Generation} & \textbf{Key Figures} & \textbf{Contribution} \\
\midrule
\textbf{Founders} (1930s--50s) & Knight, Simons, Director, Viner & Established Chicago as distinct school \\
\textbf{Core} (1950s--70s) & Friedman, Stigler, Becker & ``As if'' methodology; economic imperialism \\
\textbf{Rational Expectations} (1970s--80s) & Lucas, Sargent, Prescott & Expectations always rational \\
\textbf{Finance} (1970s--90s) & Fama, Miller, Scholes & Efficient Markets Hypothesis \\
\textbf{Law \& Economics} (1970s--) & Posner, Demsetz, Peltzman & Legal analysis as optimization \\
\bottomrule
\end{tabular}
\end{table}

\subsection{Friedman (1953): The ``F-Twist''}

Milton Friedman's essay ``The Methodology of Positive Economics'' is, according to Daniel Hausman, ``the most influential work on economic methodology of [the twentieth] century.''

The key passage:
\begin{quote}
``Truly important and significant hypotheses will be found to have `assumptions' that are wildly inaccurate descriptive representations of reality, and, in general, \textbf{the more significant the theory, the more unrealistic the assumptions}.''
\end{quote}

Paul Samuelson later termed this reasoning the ``F-Twist''---the argument that shielded assumptions (including $C_{ij} = 0$) from empirical scrutiny.

\subsection{Becker (1976): Economic Imperialism}

Gary Becker's \textit{The Economic Approach to Human Behavior} extended utility maximization to all domains:
\begin{quote}
``All human behavior can be viewed as involving participants who \textbf{maximize their utility} from a stable set of preferences.''
\end{quote}

This transformed a framework assumption into an ontological claim about human nature.

\subsection{Lucas (1976): Rational Expectations as Axiom}

Robert Lucas's ``Econometric Policy Evaluation: A Critique'' established that expectations must be model-consistent:
\begin{equation}
\mathbb{E}[C | \mathcal{I}] = C^*(\Psi) \quad \text{always}
\end{equation}

This ruled out systematic belief errors by \textit{assumption}, not evidence.

\subsection{The 2008 Crisis as Falsification}

The financial crisis exposed the limitations of treating the framework as model:

\begin{quote}
``Economists not only failed to anticipate the [financial] crisis but may have \textbf{contributed to it}---with risk and derivatives models that encouraged policy makers to see more stability than was actually present... [T]he crisis was met with incomprehension by most economists because of models that \textbf{downplay the possibility that economic actors may exhibit highly interactive behavior}.'' (Colander et al., 2009)
\end{quote}

This is precisely the consequence of treating $C_{ij} = 0$ as fact rather than simplification: models could not capture the complementarities that drove systemic risk.

\subsection{The Freshwater vs. Saltwater Divide}

The methodological divide within economics has geographic correlates:

\begin{table}[H]
\centering
\caption{Freshwater vs. Saltwater Economics}
\label{tab:freshwater-saltwater}
\renewcommand{\arraystretch}{1.3}
\small
\begin{tabular}{@{}p{5.5cm}p{5.5cm}@{}}
\toprule
\textbf{Freshwater (Chicago, Minnesota)} & \textbf{Saltwater (MIT, Harvard, Berkeley)} \\
\midrule
$C_{ij} = 0$ strictly & $C_{ij} \neq 0$ possible \\
Context exogenous, stable & Context partially endogenous \\
Markets always clear & Frictions accepted \\
Rational expectations axiom & Bounded rationality possible \\
\midrule
Friedman, Lucas, Sargent, Prescott, Fama & Samuelson, Tobin, Solow, Stiglitz, Blanchard \\
\bottomrule
\end{tabular}
\end{table}

%==============================================================================
\section{The Case for Formalization}
\label{sec:formalization}
%==============================================================================

\subsection{Why Mathematics Was Necessary}

The preceding critique of the Chicago School's epistemological confusion must not be misread as a critique of mathematical formalization itself. On the contrary: formalization was \textbf{necessary} for economics to develop as a cumulative science.

\subsection{Debreu's Defense of Rigor}

G\'{e}rard Debreu, in his 1983 Nobel Lecture and 1991 Presidential Address to the American Economic Association, articulated the core benefits:

\begin{enumerate}[nosep]
\item \textbf{Logical scrutiny:} ``In its mathematical form, economic theory is open to an efficient scrutiny for logical errors.''
\item \textbf{Cumulative development:} ``The greater logical solidity of more recent analyses has contributed to the rapid contemporary construction of economic theory. It has enabled researchers to \textbf{build on the work of their predecessors}.''
\item \textbf{Unambiguous communication:} Mathematics provides ``an efficient language for unambiguous communication among economists.''
\item \textbf{Secure foundations:} Axiomatization provides ``secure bases from which exploration could start in new dimensions.''
\end{enumerate}

\subsection{Concepts That Required Formalization}

Many of economics' most important discoveries were \textit{impossible} without mathematical formalization:

\begin{table}[H]
\centering
\caption{Concepts Requiring Mathematical Formalization}
\label{tab:formalization-required}
\renewcommand{\arraystretch}{1.3}
\small
\begin{tabular}{@{}lll@{}}
\toprule
\textbf{Concept} & \textbf{Discoverer(s)} & \textbf{Why Math Essential} \\
\midrule
Adverse Selection & Akerlof (1970) & Equilibrium with asymmetric info \\
Signaling & Spence (1973) & Separating vs. pooling equilibria \\
Moral Hazard & Mirrlees (1971) & Principal-agent optimization \\
Nash Equilibrium & Nash (1950) & Simultaneous best responses \\
Mechanism Design & Hurwicz, Maskin, Myerson & ``Reverse game theory'' \\
Arrow Impossibility & Arrow (1951) & Logical proof of impossibility \\
Option Pricing & Black-Scholes-Merton & Stochastic differential equations \\
Auction Theory & Vickrey, Myerson & Revenue equivalence theorem \\
Contract Theory & Hart, Holmstr\"{o}m & Optimal incentive structures \\
Matching Theory & Gale-Shapley, Roth & Stable matching algorithms \\
\bottomrule
\end{tabular}
\end{table}

\textbf{Without mathematical formalization}, we would not know:
\begin{itemize}[nosep]
\item Why high-quality sellers exit markets with asymmetric information
\item How to design auctions that reveal true valuations
\item When additional performance signals improve contracts
\item That no voting system can satisfy all fairness criteria simultaneously
\end{itemize}

\subsection{Rigor vs. ``Mathiness''}

Paul Romer (Nobel 2018) distinguished proper mathematical rigor from its abuse:
\begin{quote}
``Mathiness... is a style that lets economists draw \textbf{invalid inferences} from the assumptions and structure of a model.''
\end{quote}

The distinction:
\begin{table}[H]
\centering
\renewcommand{\arraystretch}{1.3}
\begin{tabular}{@{}p{5.5cm}p{5.5cm}@{}}
\toprule
\textbf{Rigor (Debreu-Style)} & \textbf{Mathiness (Abuse)} \\
\midrule
Mathematics makes assumptions explicit & Mathematics obscures assumptions \\
Conclusions follow logically & Invalid inferences from models \\
Limitations acknowledged & Limitations ignored \\
Framework $\neq$ Model & Framework $=$ Model \\
\bottomrule
\end{tabular}
\end{table}

\subsection{The Real Critique}

The critique of the Chicago School is \textbf{not}:
\begin{quote}
``Economics should have used less mathematics.''
\end{quote}

The critique \textbf{is}:
\begin{quote}
``Economics should have maintained the epistemological distinction between framework and model, treating $C_{ij} = 0$ as a \textit{simplifying assumption} rather than an \textit{empirical claim}.''
\end{quote}

The mathematics was correct. The epistemological confusion was the error.

%==============================================================================
\section{Comparative Perspective: The Alternative to Formalization}
\label{sec:comparative}
%==============================================================================

To understand why formalization was necessary, consider what happened in social sciences that \textit{did not} formalize.

\subsection{Psychology: The Replication Crisis}

Psychology, which relied less on mathematical formalization, experienced a devastating replication crisis:
\begin{itemize}
\item Only 36\% of 100 studies replicated (Open Science Collaboration, 2015)
\item ``P-hacking'' widespread (Simmons, Nelson, Simonsohn, 2011)
\item Publication bias severe (Rosenthal, 1979: ``file drawer problem'')
\item Major effects not robust: ego depletion, power posing, social priming
\end{itemize}

\textbf{Without mathematical rigor}: Decades of research on ``ego depletion,'' ``power posing,'' and ``social priming'' proved largely non-replicable. Verbal theories permitted ambiguity that formal models would have exposed.

\subsection{Sociology: Paradigm Fragmentation}

Sociology, explicitly rejecting mathematical formalization, exhibits:
\begin{itemize}[nosep]
\item Multiple competing paradigms with no integration mechanism
\item Persistent debates over ``What did Weber \textit{really} mean?''
\item Limited policy influence compared to economics
\item No cumulative theoretical development
\end{itemize}

As Randall Collins observed: ``Sociology remains a discipline with many paradigms but no consensus.''

\subsection{Management Science: Fashion Cycles}

Business and management research, with minimal mathematical formalization, exhibits fashion cycles rather than cumulative progress:
\begin{center}
1980s: TQM $\to$ 1990s: Reengineering $\to$ 2000s: Six Sigma $\to$ 2010s: Agile $\to$ 2020s: Digital Transformation
\end{center}

Each generation ``forgets'' the previous. No cumulative knowledge base develops. Consultants sell ``new'' concepts that recycle old ideas.

\subsection{The Comparative Verdict}

\begin{table}[H]
\centering
\caption{Formalized vs. Non-Formalized Social Sciences}
\label{tab:comparative}
\renewcommand{\arraystretch}{1.3}
\small
\begin{tabular}{@{}lcccc@{}}
\toprule
\textbf{Criterion} & \textbf{Economics} & \textbf{Psychology} & \textbf{Sociology} & \textbf{Management} \\
\midrule
Mathematical formalization & High & Low & Very low & Very low \\
Clear definitions & Yes & Partial & No & No \\
Cumulative development & Yes & Weak & No & No \\
Replicability & High & 36\% & Unknown & N/A \\
Policy influence & Very high & Moderate & Low & High (fashion) \\
Predictive power & Moderate & Weak & Very weak & None \\
Nobel Prize & Since 1969 & No & No & No \\
\bottomrule
\end{tabular}
\end{table}

\textbf{The pattern is clear:} Formalization correlates with cumulation and influence. Its absence correlates with fragmentation and marginalization.

%==============================================================================
\section{Technological Preconditions}
\label{sec:technology}
%==============================================================================

\subsection{The Computational Constraint}

The preceding analysis might suggest that the Chicago School made an avoidable epistemological error. This would be historically naive. The framework-model conflation was not merely intellectually convenient---it was \textbf{technologically inevitable} given the computational constraints of the era.

\begin{table}[H]
\centering
\caption{Computational Technology and Model Complexity}
\label{tab:tech-constraints}
\renewcommand{\arraystretch}{1.3}
\small
\begin{tabular}{@{}lll@{}}
\toprule
\textbf{Era} & \textbf{Technology} & \textbf{Modeling Constraint} \\
\midrule
1940s--1960s & Pencil, slide rule, early mainframes & Only analytically solvable models \\
1970s--1980s & Mainframes, early PCs & Simple numerical solutions possible \\
1990s--2000s & PCs, early internet & Monte Carlo, agent-based models emerge \\
2010s & Cloud computing, big data & Large-scale simulations feasible \\
2020s & AI/LLMs, massive compute & Context-dependent models operationalizable \\
\bottomrule
\end{tabular}
\end{table}

\subsection{Why $C_{ij} = 0$ Was Necessary}

The separability assumption ($C_{ij} = 0$) was not chosen because economists believed it was true. It was chosen because it was \textbf{the only assumption that permitted analytical solutions}:

\begin{table}[H]
\centering
\renewcommand{\arraystretch}{1.3}
\begin{tabular}{@{}p{5.5cm}p{5.5cm}@{}}
\toprule
\textbf{With $C_{ij} = 0$} & \textbf{With $C_{ij} \neq 0$} \\
\midrule
Utility separable across goods & Cross-effects everywhere \\
Individual optimization independent & Simultaneous optimization required \\
Equilibrium: solve $n$ equations & Equilibrium: solve $n^2$ interactions \\
Analytical solutions possible & Numerical methods required \\
Pencil and paper sufficient & Computers required \\
\bottomrule
\end{tabular}
\end{table}

\subsection{Friedman's Defense Revisited}

Friedman's ``as if'' methodology (1953) can be read not as epistemological confusion but as \textbf{pragmatic necessity}:

\begin{quote}
\textit{Given that we cannot compute models with realistic assumptions, we must use unrealistic assumptions that permit computation. The test is predictive success, not assumption realism.}
\end{quote}

This was \textbf{defensible in 1953}. It became problematic only when the technological constraints lifted but the methodology persisted.

\subsection{The Technological Milestones}

\begin{table}[H]
\centering
\caption{Technology Enabling Framework-Model Separation}
\label{tab:tech-milestones}
\renewcommand{\arraystretch}{1.3}
\small
\begin{tabular}{@{}llp{6cm}@{}}
\toprule
\textbf{Year} & \textbf{Technology} & \textbf{What It Enabled} \\
\midrule
1981 & IBM PC & Econometrics democratized \\
1985 & Spreadsheets & First simulations accessible \\
1991 & World Wide Web & Data sharing, collaboration \\
2000s & Cloud computing & Massive Monte Carlo simulations \\
2007 & Amazon MTurk & Cheap experimental subjects \\
2008 & Smartphones & Real-time behavioral data \\
2017 & Transformer architecture & Language models feasible \\
2022 & ChatGPT/Claude & LLM Monte Carlo for context estimation \\
\bottomrule
\end{tabular}
\end{table}

\subsection{What AI Specifically Enables}

The final technological piece---large language models---enables what was previously impossible: \textbf{operationalizing context-dependent models at scale}.

\begin{table}[H]
\centering
\caption{AI-Enabled Capabilities for Integrative Frameworks}
\label{tab:ai-enables}
\renewcommand{\arraystretch}{1.3}
\small
\begin{tabular}{@{}lp{4.5cm}p{4.5cm}@{}}
\toprule
\textbf{Component} & \textbf{Before AI} & \textbf{With AI} \\
\midrule
Context estimation & Manual surveys, expensive, slow & LLM Monte Carlo simulation \\
Context-specific models & Develop each separately & Generate automatically from framework \\
Heterogeneous preferences & 2--3 discrete types & Continuous distributions \\
Cross-domain integration & Manual, inconsistent & Embedding-based consistency \\
Calibration & Months of fieldwork & Hours with LLM simulation \\
\bottomrule
\end{tabular}
\end{table}

\subsection{The Historical Verdict}

\textbf{Retrospectively:} The framework-model conflation was epistemologically incorrect. Treating $C_{ij} = 0$ as empirical fact rather than simplifying assumption led to 50+ years of fragmentation.

\textbf{Practically:} The conflation was defensible given technological constraints. Without computers, models with $C_{ij} \neq 0$ were unsolvable. The Chicago methodology was a rational response to computational limitations.

\textbf{Now:} The technological constraints have lifted. Computers can solve models with $C_{ij} \neq 0$. AI can estimate context across situations. The framework-model distinction is now \textit{operationally feasible}.

%==============================================================================
\section{Integration Outside Academia: A Structural Prediction}
\label{sec:outside}
%==============================================================================

\subsection{The Prediction}

If the barrier to integration is structural, integration should emerge from contexts with different structures. This prediction is consistent with observed patterns:

\begin{itemize}[nosep]
\item Consulting firms integrate models for practical application
\item Policy organizations need unified frameworks for decision-making
\item Research institutes outside universities face different incentives
\end{itemize}

\subsection{Why Integration Succeeds Outside Academia}

\begin{enumerate}
\item \textbf{Different Incentives:} Consulting firms need working models, not publications
\item \textbf{Practice Pressure:} Transformation projects fail at 70\%---integration is demanded
\item \textbf{Long Time Horizon:} 20+ years of iteration, not 3-year grant cycles
\item \textbf{Interdisciplinary Teams:} Economists + psychologists + practitioners
\item \textbf{No Territorial Stake:} No paradigm to defend
\end{enumerate}

\subsection{The Integration Proof}

\textbf{Claim:} Integration is possible.

\textbf{Evidence:}
\begin{itemize}
\item Integrative frameworks exist and are used in practice
\item Multiple disciplinary insights can be unified under common concepts
\item Existing models become \textit{special cases}, not competitors
\item Framework + calibration data = predictive power
\end{itemize}

\textbf{Conclusion:} The barrier was institutional, not intellectual. The framework-model confusion prevented integration for 50+ years.

\subsection{How Integrative Frameworks Address Each Barrier}

\begin{table}[H]
\centering
\caption{Addressing Structural Barriers}
\label{tab:ebf-response}
\renewcommand{\arraystretch}{1.4}
\begin{tabular}{@{}p{4cm}p{8cm}@{}}
\toprule
\textbf{Barrier} & \textbf{Response} \\
\midrule
Academic incentives & Built outside academia; validated through practice \\
Disciplinary silos & Single conceptual framework spans all disciplines \\
Peer review bias & Validation through application, not publication \\
Funding structures & Funded by practice needs, not grants \\
Territorial defense & Shows models as \textit{special cases}, preserving credit \\
Replication crisis & Integration enables cross-model consistency checks \\
\bottomrule
\end{tabular}
\end{table}

%==============================================================================
\section{The Metascience Movement}
\label{sec:metascience}
%==============================================================================

\subsection{Key Institutions}

\begin{itemize}
\item \textbf{Center for Open Science} (Brian Nosek): Pre-registration, open data
\item \textbf{METRICS} (Stanford, Ioannidis): Meta-research on research
\item \textbf{Research on Research Institute} (UK): Systemic analysis
\item \textbf{Santa Fe Institute}: Complexity as integration principle
\end{itemize}

\subsection{Key Reforms Proposed}

\begin{enumerate}
\item \textbf{Pre-registration:} Commit to hypotheses before data collection
\item \textbf{Registered Reports:} Peer review before results known
\item \textbf{Open Data:} Share data for replication
\item \textbf{Multi-site Studies:} Replicate across contexts
\item \textbf{Meta-analysis:} Cumulative synthesis
\end{enumerate}

\subsection{What's Missing: Integration}

The metascience movement focuses on \textit{replication} and \textit{transparency}---but not on \textit{integration across paradigms}.

Integration fills this gap: It doesn't just replicate existing models, it \textbf{unifies} them under common organizing principles.

%==============================================================================
\section{Economics as the Formalizable Social Science}
\label{sec:thesis}
%==============================================================================

\subsection{The Central Thesis}

Economics is the \textbf{only fully formalizable social science}---not because economists were smarter, but because they \textit{pragmatically adapted} to technological constraints.

The adaptation produced epistemological confusion. But it also produced \textit{the infrastructure} for a complete science of human behavior.

\subsection{The Contrast with Other Disciplines}

\begin{table}[H]
\centering
\caption{Economics vs. Other Social Sciences: The Formalizability Thesis}
\label{tab:formalizability}
\renewcommand{\arraystretch}{1.3}
\small
\begin{tabular}{@{}p{3cm}p{4.5cm}p{4.5cm}@{}}
\toprule
\textbf{Discipline} & \textbf{Approach to Technology} & \textbf{Result} \\
\midrule
\textbf{Economics} & Adapted methods to what was computable; formalized everything possible & Fully formalizable; cumulative; policy-relevant \\
\textbf{Psychology} & Rejected formalization as ``reductionist'' & Replication crisis; limited cumulation \\
\textbf{Sociology} & Explicitly rejected mathematical methods & Paradigm fragmentation; no integration \\
\textbf{Political Science} & Partial formalization (rational choice) & Split discipline; limited prediction \\
\textbf{Anthropology} & Rejected quantification entirely & Rich description; no generalization \\
\bottomrule
\end{tabular}
\end{table}

\subsection{The Infrastructure Argument}

Seventy years of economic formalization created something irreplaceable: a \textbf{mathematical backbone} onto which insights from all social sciences can be integrated.

Integration can only emerge from the economic tradition \textbf{not} because:
\begin{itemize}[nosep]
\item Economists are smarter (they are not)
\item Economic insights are more important (all disciplines contribute)
\item Economics ``colonized'' other fields (imperialism is not integration)
\end{itemize}

Integration can only emerge from economics \textbf{because}:
\begin{itemize}[nosep]
\item Economics built the formal infrastructure over 70 years
\item This infrastructure provides the backbone for integration
\item Other social sciences lack this cumulative formal structure
\item AI can operationalize this infrastructure---but cannot create it
\end{itemize}

\textbf{The timing matters:} The infrastructure had to exist \emph{before} AI could operationalize it.

%==============================================================================
\section{Discussion}
\label{sec:discussion}
%==============================================================================

\subsection{The Core Paradox Revisited}

The central paradox of behavioral science fragmentation can be stated simply:

\begin{quote}
\textit{Everyone agrees integration is needed. No one is rewarded for doing it.}
\end{quote}

This is not a coordination failure that can be solved by agreement. It is a structural problem requiring structural solutions.

\subsection{Limitations}

This analysis has limitations:
\begin{itemize}[nosep]
\item Evidence is largely observational; causal identification is difficult
\item Quantitative estimates vary across studies and contexts
\item Counter-examples exist (successful interdisciplinary fields)
\item The analysis focuses on barriers, not detailed solutions
\end{itemize}

\subsection{Future Research}

Promising directions include:
\begin{itemize}[nosep]
\item Experimental studies of peer review for integrative work
\item Analysis of successful integration cases
\item Design of institutional mechanisms that reward integration
\item Evaluation of alternative metrics for scientific contribution
\item Longitudinal studies of framework adoption in practice
\end{itemize}

\subsection{Institutional Innovation}

Overcoming structural barriers requires institutional innovation:

\textbf{Alternative incentive structures:} Organizations that reward integration directly, not through publication metrics.

\textbf{Boundary-spanning roles:} Scientists who specialize in translation across disciplines, valued for integration rather than within-discipline contribution.

\textbf{Different funding mechanisms:} Support for long-term, high-risk integration projects with appropriate evaluation criteria.

\textbf{New publication venues:} Journals that value integration, with review processes appropriate for interdisciplinary work.

%==============================================================================
\section{Conclusion}
\label{sec:conclusion}
%==============================================================================

The fragmentation of behavioral science into 40+ competing models is not an intellectual failure but a structural outcome. Five mechanisms---academic incentives, disciplinary silos, peer review bias, funding structures, and territorial defense---systematically prevent integration despite universal recognition of its value.

A central epistemological error---treating Samuelson's framework as a falsifiable model---contributed to fragmentation when behavioral anomalies accumulated without an integrative structure to absorb them. This error was technologically defensible when only analytically solvable models were computationally feasible, but became problematic as computational constraints lifted.

The barrier to integration is institutional, not intellectual. Integration is technically possible---the existence of integrative frameworks demonstrates this. But academic institutions as currently structured will not produce it.

This analysis suggests that integration will continue to emerge from outside traditional academic structures: consulting firms, policy organizations, and research institutes with different incentive systems. The Evidence-Based Framework for Economic and Social Behavior exemplifies this pattern: built outside academia, validated through practice, and designed to unify rather than compete.

For behavioral science to achieve cumulative integration, institutional innovation is required. The science of science has documented the problem; designing solutions remains for future work. The technological preconditions for integration now exist---what remains is the institutional will to pursue it.

%==============================================================================
% References
%==============================================================================
\bibliographystyle{apalike}

\begin{thebibliography}{99}

\bibitem[Akerlof, 1970]{akerlof1970}
Akerlof, G. A. (1970). The market for ``lemons'': Quality uncertainty and the market mechanism. \textit{Quarterly Journal of Economics}, 84(3), 488--500.

\bibitem[Arrow, 1951]{arrow1951}
Arrow, K. J. (1951). Social choice and individual values. \textit{Cowles Foundation Monographs}, No. 12. New York: Wiley.

\bibitem[Backhouse \& Cherrier, 2017]{backhouse2017}
Backhouse, R. E., \& Cherrier, B. (2017). The age of the applied economist: The transformation of economics since the 1970s. \textit{History of Political Economy}, 49(Supplement), 1--33.

\bibitem[Bammer, 2013]{bammer2013}
Bammer, G. (2013). \textit{Disciplining Interdisciplinarity: Integration and Implementation Sciences for Researching Complex Real-World Problems}. Canberra: ANU Press.

\bibitem[Becker, 1976]{becker1976}
Becker, G. S. (1976). \textit{The Economic Approach to Human Behavior}. Chicago: University of Chicago Press.

\bibitem[Blaug, 1992]{blaug1992}
Blaug, M. (1992). \textit{The Methodology of Economics: Or, How Economists Explain} (2nd ed.). Cambridge: Cambridge University Press.

\bibitem[Bloor, 1976]{bloor1976}
Bloor, D. (1976). \textit{Knowledge and Social Imagery}. London: Routledge \& Kegan Paul.

\bibitem[Burgin, 2012]{burgin2012}
Burgin, A. (2012). \textit{The Great Persuasion: Reinventing Free Markets since the Depression}. Cambridge, MA: Harvard University Press.

\bibitem[Caldwell, 1980]{caldwell1980}
Caldwell, B. J. (1980). A critique of Friedman's methodological instrumentalism. \textit{Southern Economic Journal}, 47(2), 366--374.

\bibitem[Camerer et al., 2016]{camerer2016}
Camerer, C. F., et al. (2016). Evaluating replicability of laboratory experiments in economics. \textit{Science}, 351(6280), 1433--1436.

\bibitem[Camerer et al., 2018]{camerer2018}
Camerer, C. F., et al. (2018). Evaluating the replicability of social science experiments in Nature and Science between 2010 and 2015. \textit{Nature Human Behaviour}, 2(9), 637--644.

\bibitem[Colander, 2004]{colander2004changing}
Colander, D. (2004). The changing face of mainstream economics. \textit{Review of Political Economy}, 16(4), 485--499.

\bibitem[Colander et al., 2009]{colander2009crisis}
Colander, D., et al. (2009). The financial crisis and the systemic failure of the economics profession. \textit{Critical Review}, 21(2--3), 249--267.

\bibitem[Davis, 2008]{davis2008heterodox}
Davis, J. B. (2008). The turn in recent economics and return of orthodoxy. \textit{Cambridge Journal of Economics}, 32(3), 349--366.

\bibitem[Debreu, 1984]{debreu1984}
Debreu, G. (1984). Economic theory in the mathematical mode. \textit{American Economic Review}, 74(3), 267--278.

\bibitem[Debreu, 1991]{debreu1991}
Debreu, G. (1991). The mathematization of economic theory. \textit{American Economic Review}, 81(1), 1--7.

\bibitem[Edwards \& Roy, 2017]{edwards2017}
Edwards, M. A., \& Roy, S. (2017). Academic research in the 21st century: Maintaining scientific integrity in a climate of perverse incentives and hypercompetition. \textit{Environmental Engineering Science}, 34(1), 51--61.

\bibitem[Fortunato et al., 2018]{fortunato2018}
Fortunato, S., et al. (2018). Science of science. \textit{Science}, 359(6379), eaao0185.

\bibitem[Fourcade, 2009]{fourcade2009}
Fourcade, M. (2009). \textit{Economists and Societies: Discipline and Profession in the United States, Britain, and France, 1890s to 1990s}. Princeton: Princeton University Press.

\bibitem[Fourcade et al., 2015]{fourcade2015superiority}
Fourcade, M., Ollion, E., \& Algan, Y. (2015). The superiority of economists. \textit{Journal of Economic Perspectives}, 29(1), 89--114.

\bibitem[Friedman, 1953]{friedman1953}
Friedman, M. (1953). The methodology of positive economics. In \textit{Essays in Positive Economics} (pp. 3--43). Chicago: University of Chicago Press.

\bibitem[Gintis, 2009]{gintis2009}
Gintis, H. (2009). \textit{The Bounds of Reason: Game Theory and the Unification of the Behavioral Sciences}. Princeton: Princeton University Press.

\bibitem[Hands, 2001]{hands2001}
Hands, D. W. (2001). \textit{Reflection Without Rules: Economic Methodology and Contemporary Science Theory}. Cambridge: Cambridge University Press.

\bibitem[Hausman, 1992]{hausman1992}
Hausman, D. M. (1992). \textit{The Inexact and Separate Science of Economics}. Cambridge: Cambridge University Press.

\bibitem[Ioannidis, 2005]{PAP-ioannidis2005}
Ioannidis, J. P. A. (2005). Why most published research findings are false. \textit{PLoS Medicine}, 2(8), e124.

\bibitem[Jacobs, 2013]{jacobs2013}
Jacobs, J. A. (2013). \textit{In Defense of Disciplines: Interdisciplinarity and Specialization in the Research University}. Chicago: University of Chicago Press.

\bibitem[Klein, 1990]{klein1990}
Klein, J. T. (1990). \textit{Interdisciplinarity: History, Theory, and Practice}. Detroit: Wayne State University Press.

\bibitem[Koopmans, 1957]{koopmans1957}
Koopmans, T. C. (1957). \textit{Three Essays on the State of Economic Science}. New York: McGraw-Hill.

\bibitem[Kuhn, 1962]{kuhn1962}
Kuhn, T. S. (1962). \textit{The Structure of Scientific Revolutions}. Chicago: University of Chicago Press.

\bibitem[Lakatos, 1970]{lakatos1970}
Lakatos, I. (1970). Falsification and the methodology of scientific research programmes. In I. Lakatos \& A. Musgrave (Eds.), \textit{Criticism and the Growth of Knowledge} (pp. 91--196). Cambridge: Cambridge University Press.

\bibitem[Latour \& Woolgar, 1979]{latour1979}
Latour, B., \& Woolgar, S. (1979). \textit{Laboratory Life: The Construction of Scientific Facts}. Beverly Hills: Sage.

\bibitem[Lawson, 2006]{lawson2006nature}
Lawson, T. (2006). The nature of heterodox economics. \textit{Cambridge Journal of Economics}, 30(4), 483--505.

\bibitem[Lazear, 2000]{lazear2000}
Lazear, E. P. (2000). Economic imperialism. \textit{Quarterly Journal of Economics}, 115(1), 99--146.

\bibitem[Lucas, 1976]{lucas1976}
Lucas, R. E., Jr. (1976). Econometric policy evaluation: A critique. \textit{Carnegie-Rochester Conference Series on Public Policy}, 1, 19--46.

\bibitem[Medema, 2023]{medema2023chicago}
Medema, S. G. (2023). Identifying a ``Chicago School'' of economics: On the origins, diffusion, and evolving meanings of a famous brand name. \textit{History of Political Economy}, 55(S1), 51--74.

\bibitem[Merton, 1942]{merton1942}
Merton, R. K. (1942). The normative structure of science. In \textit{The Sociology of Science} (pp. 267--278). Chicago: University of Chicago Press.

\bibitem[Open Science Collaboration, 2015]{osc2015}
Open Science Collaboration. (2015). Estimating the reproducibility of psychological science. \textit{Science}, 349(6251), aac4716.

\bibitem[Popper, 1959]{popper1959}
Popper, K. R. (1959). \textit{The Logic of Scientific Discovery}. London: Hutchinson.

\bibitem[Reder, 1982]{reder1982}
Reder, M. W. (1982). Chicago economics: Permanence and change. \textit{Journal of Economic Literature}, 20(1), 1--38.

\bibitem[Romer, 2015]{romer2015}
Romer, P. M. (2015). Mathiness in the theory of economic growth. \textit{American Economic Review}, 105(5), 89--93.

\bibitem[Samuelson, 1947]{samuelson1947}
Samuelson, P. A. (1947). \textit{Foundations of Economic Analysis}. Cambridge, MA: Harvard University Press.

\bibitem[Samuelson, 1963]{samuelson1963}
Samuelson, P. A. (1963). Problems of methodology: Discussion. \textit{American Economic Review}, 53(2), 231--236.

\bibitem[Smaldino \& McElreath, 2016]{smaldino2016}
Smaldino, P. E., \& McElreath, R. (2016). The natural selection of bad science. \textit{Royal Society Open Science}, 3(9), 160384.

\bibitem[Stephan, 2012]{stephan2012}
Stephan, P. E. (2012). \textit{How Economics Shapes Science}. Cambridge, MA: Harvard University Press.

\bibitem[Uzzi et al., 2013]{uzzi2013}
Uzzi, B., Mukherjee, S., Stringer, M., \& Jones, B. (2013). Atypical combinations and scientific impact. \textit{Science}, 342(6157), 468--472.

\bibitem[Van Horn \& Mirowski, 2009]{vanhorn2009}
Van Horn, R., \& Mirowski, P. (2009). The rise of the Chicago School of Economics and the birth of neoliberalism. In P. Mirowski \& D. Plehwe (Eds.), \textit{The Road from Mont P\`{e}lerin} (pp. 139--178). Cambridge, MA: Harvard University Press.

\bibitem[Van Horn et al., 2011]{vanhorn2011}
Van Horn, R., Mirowski, P., \& Stapleford, T. A. (Eds.). (2011). \textit{Building Chicago Economics: New Perspectives on the History of America's Most Powerful Economics Program}. Cambridge: Cambridge University Press.

\bibitem[Wang et al., 2017]{wang2017}
Wang, J., Veugelers, R., \& Stephan, P. (2017). Bias against novelty in science: A cautionary tale for users of bibliometric indicators. \textit{Research Policy}, 46(8), 1416--1436.

\bibitem[Wilson, 1998]{wilson1998}
Wilson, E. O. (1998). \textit{Consilience: The Unity of Knowledge}. New York: Knopf.

\bibitem[Yonay, 1998]{yonay1998}
Yonay, Y. P. (1998). \textit{The Struggle over the Soul of Economics: Institutionalist and Neoclassical Economists in America between the Wars}. Princeton: Princeton University Press.

\end{thebibliography}

%==============================================================================
% Appendix
%==============================================================================
\appendix
\section{Data Sources}
\label{app:data}

This paper synthesizes findings from the following primary data sources:
\begin{itemize}[nosep]
\item Web of Science citation database (Fortunato et al., 2018)
\item NIH and NSF grant databases (Stephan, 2012)
\item Peer review outcome data (Wang et al., 2017)
\item Agent-based simulations (Smaldino \& McElreath, 2016)
\item Open Science Collaboration replication data (2015)
\item Economics replication data (Camerer et al., 2016, 2018)
\end{itemize}

\section{Historical Timeline}
\label{app:timeline}

\begin{table}[H]
\centering
\caption{Key Events in the Framework-Model Confusion}
\renewcommand{\arraystretch}{1.2}
\small
\begin{tabular}{@{}ll@{}}
\toprule
\textbf{Year} & \textbf{Event} \\
\midrule
1947 & Samuelson publishes \textit{Foundations of Economic Analysis} \\
1953 & Friedman publishes ``Methodology of Positive Economics'' \\
1959 & Arrow-Debreu prove general equilibrium existence \\
1962 & Kuhn publishes \textit{Structure of Scientific Revolutions} \\
1970 & Akerlof publishes ``Market for Lemons'' \\
1976 & Lucas Critique published \\
1979 & Kahneman-Tversky publish Prospect Theory \\
1999 & Fehr-Schmidt publish inequality aversion model \\
2005 & Ioannidis publishes ``Why Most Research Findings Are False'' \\
2008 & Financial crisis exposes model limitations \\
2015 & Open Science Collaboration reports 36\% replication rate \\
2016 & Smaldino \& McElreath publish ``Natural Selection of Bad Science'' \\
2020s & AI enables context-dependent modeling at scale \\
\bottomrule
\end{tabular}
\end{table}

\section{Relationship to EBF}
\label{app:ebf}

This working paper was generated from Appendix LIT-META of the Evidence-Based Framework documentation. The appendix contains additional detail on:
\begin{itemize}[nosep]
\item The 8$\Psi$ Model of Scientific Development
\item Beatrix operationalization of the framework-model distinction
\item Extended historical case studies
\item Quantitative citation network analysis
\item EBF-specific responses to metascience critique
\end{itemize}

For the complete analysis, see Appendix LIT-META in the EBF documentation.

\end{document}
